% chapters/12/sub-cpmk.tex
\begin{subcpmkbox}
Bab ini membuka blok baru, yaitu \CPMK{3}. Perkakasnya berganti dari pustaka
menjadi mesin permainan, sedangkan konsepnya sebagian besar telah Anda kenal.
Bab ini menuntaskan \SUBCPMK{3.1} secara utuh.

\vspace{0.5em}
\textbf{\SUBCPMK{3.1} Menavigasi antarmuka Unity dan menyusun
\istilah{project}, \istilah{scene}, serta \istilah{GameObject} dasar} ---
ranah C3, \CPMK{3}, \CPL{04}.

\vspace{0.5em}
{\setlength{\tabcolsep}{4pt}\footnotesize%
\begin{tabularx}{\textwidth}{@{}C{1.9cm} L{4.0cm} L{2.5cm} Y@{}}
\toprule
\headrow \textbf{Indikator} & \textbf{Rumusan} & \textbf{Dibahas pada} & \textbf{Bukti} \\
\midrule
\ind{a} & Mengoperasikan \istilah{Hierarchy}, \istilah{Inspector},
\istilah{Project window}, dan \istilah{Scene} &
Bagian~\ref{sec:antarmuka-unity} &
Catatan praktikum yang menyebutkan kegunaan setiap panel, dan tangkapan layar
scene buatan sendiri. \\
\addlinespace
\ind{b} & Membuat dan menyimpan \istilah{scene} berisi kamera serta objek
primitif &
Bagian~\ref{sec:scene-project} dan praktikum pada akhir bab &
Berkas \istilah{scene} yang tersimpan, berisi kamera, pencahayaan, dan objek
primitif beserta \istilah{script} C\# yang menempel padanya. \\
\addlinespace
\ind{c} & Menjelaskan padanan konsep Unity dengan konsep Pygame yang telah
dipelajari &
Bagian~\ref{sec:padanan-unity} &
Tabel padanan yang disusun sendiri, beserta penjelasan singkat setiap
pasangannya. \\
\bottomrule
\end{tabularx}}
\end{subcpmkbox}
